%------------------------------------------------------------------------------
% C-- specific shortcuts
%------------------------------------------------------------------------------
% claude: the original Quick C-- documentation defined these in doc/macros.tex
% (kept here as ../macros.tex). We do not \input that file because its second
% half redefines the noweb \nwstartdeflinemarkup/\nwenddeflinemarkup hooks and
% does a \RequirePackage{noweb}, which conflicts with the latex/noweb.sty
% loaded from Cminusminus.nw. So the shorthands are redefined here in the
% principia style (see latex/Macros.tex and its \plan / \unix).

% claude: the "Extra Code" appendix holds one chapter per source directory, so
% the book has more than 26 appendix chapters and the default \Alph numbering
% overflows with "Counter too large". \AlphAlph continues with AA, AB, ...
\usepackage{alphalph}

\newcommand{\PAL}{C-{}-\xspace}
\let\C\PAL
\let\pal\PAL
\newcommand{\qcc}{QC-{}-\xspace}

\newcommand{\asdl}{ASDL\xspace}
\newcommand{\rtl}{RTL\xspace}
\newcommand{\ir}{IR\xspace}
\newcommand{\AST}{AST\xspace}
\newcommand{\cfg}{CFG\xspace}
\newcommand{\burg}{BURG\xspace}
\newcommand{\mips}{MIPS\xspace}
\newcommand{\sled}{SLED\xspace}
\newcommand{\sparc}{SPARC\xspace}

% claude: used in math mode too, hence the \ifmmode dance (from ../macros.tex)
\def\vfp{\relax\ifmmode\mathit{VFP}\else VFP\fi}
\def\sp{\relax\ifmmode\mathit{SP}\else SP\fi}
% claude: the C-- papers write assignment as ':=', not as the LaTeX arrow
\def\gets{\mathrel{:=}}

\newcommand{\module}[1]{\texttt{#1.nw}}
\def\remark#1{\marginpar{\raggedright\hbadness=10000
        \def\baselinestretch{0.8}\scriptsize
        \it #1\par}}
\def\interface#1{\section{Interface: #1}}
\def\implementation#1{\section{Implementation: #1}}
\newenvironment{cl}{\begin{quote}\it}{\end{quote}}
